\section{The horizontal middle: width and curvature estimates}
\label{app:horizontal-estimates}

Throughout this appendix $L=\pi/2$. For a centered height-one cap
write $I=[-2C,2C]$, $J=[-C,C]$, $A=f_U$ and $n=n_L(U)$.
A horizontal canonical middle has $A=1$ on $J$, since the top
face meets $J$ by \cref{lem:selection-top}. We prove the following
input to the full-turn argument.

\begin{lemma}[Horizontal unit curvature]\label{lem:full-horizontal-unit}
Every canonical full maximum with horizontal middle satisfies
\[
 \frac12<C<C_H,\qquad
 C_H=\frac1{35}\sqrt{523+2\sqrt{701}}.
\]
Its top face is exactly $J\times\{1\}$, and its two regular
curvature densities satisfy $0\le u,v\le1$ almost everywhere.
There are unique $0<\beta<\gamma<L$ such that
\[
 \begin{array}{lll}
 p>0&(0,\beta),&p<0\quad(\beta,L),\\
 q>0&(0,\gamma),&q<0\quad(\gamma,L).
 \end{array}
\]
\end{lemma}

The proof uses the endpoint and source laws of
\cref{sec:variations}; the scalar exclusions below use ordinary
areas of actual cap and niche pieces. Put $k=\sqrt2-1$.

\subsection{A three-angle exclusion of short widths}

For this subsection use $a=2C$ and write the actual diagonal
support data as
\[
 \begin{gathered}
 \xi=\sqrt2h_U(\pi/4)-1,\qquad
 \zeta=\sqrt2h_U(3\pi/4)-1,\\
 m=(\xi+\zeta)/2,\quad d=(\xi-\zeta)/2.
 \end{gathered}
\]
The plateau and containing box imply $a/2\le\xi,\zeta\le a$.
The roof is bounded by
$A_0(x)=\min(1,1+\xi-x,1+\zeta+x)$, whereas the angle
$\pi/4$ puts the tent
$n_{45}=[\min(\xi-k-x,\zeta-k+x)]_+$ below $n$.
Its apex $d$ lies in $J$. For $a>k$, integration gives
\begin{equation}\label{eq:horizontal-diagonal}
 \begin{split}
 P_L(U)&\le\Phi(a,\xi,\zeta)-D,\\
 \Phi&=a-\frac{(a-\xi)^2+(a-\zeta)^2}{2}-(m-k)_+^2\\
 &\hspace{9mm}
 +\frac{(\xi-k-a/2)_+^2+(\zeta-k-a/2)_+^2}{2},\\
 D&=\int_{I\setminus J}(A_0-A)\ge0.
 \end{split}
\end{equation}
For fixed $a$, $\Phi$ is symmetric and concave in $(\xi,\zeta)$.
Where $m>k$ its Hessian is
\[
 \begin{pmatrix}
 -3/2+\mathbf1_{\xi>k+a/2}&-1/2\\
 -1/2&-3/2+\mathbf1_{\zeta>k+a/2}
 \end{pmatrix};
\]
each of the four matrices is negative semidefinite. Where $m\le k$
the clipping terms vanish and the Hessian is $-\mathrm{Id}$.
The first derivatives agree across the boundaries. In particular
$\Phi(a,\xi,\zeta)\le\Phi(a,m,m)$.

\begin{lemma}[Short horizontal caps]\label{lem:horizontal-short}
If $0<a\le1$ and $A=1$ on $[-a/2,a/2]$, then
$P_L(U)\le41/50<\Mc/2$.
\end{lemma}
\begin{proof}
For $a\le k$ the bound $P_L(U)\le a$ suffices. Otherwise the
symmetric critical point $\xi=\zeta=(a+k)/2$ of \eqref{eq:horizontal-diagonal}
is a concave maximum, giving
\[
 \Phi\le B(a):=a-(a-k)^2/2\le B(1)=2k.
\]
Suppose $P_L(U)>41/50$. The following restrictions follow from
this same concave function:
\begin{equation}\label{eq:horizontal-short-data}
 a>9/10,\quad m>16/25,\quad
 \xi-a/2,\ \zeta-a/2<8/25,\quad D<3/350.
\end{equation}
Here are the exact optimizations. $B$ increases up to $a=1$,
and $B(9/10)<9/10-(2/5)^2/2=41/50$ because $k<1/2$.
For $m\le16/25$ the symmetric clipping terms vanish.
If $m\le1/2$ then
$a-(a-m)^2=m+1/4-(a-m-1/2)^2\le3/4$.
For $1/2<m\le16/25$ the expression increases in $a\le1$ and
then in $m\le16/25$, with upper value
$(2050k-337)/625<41/50$.
On the constrained half-square $\zeta-a/2\ge t_0=8/25$,
the concave maximum is at
\[
 \zeta=a/2+t_0,\qquad \xi=a/2+(2k-t_0)/3.
\]
Its clipping terms vanish, its $\xi$ derivative is zero and its
$\zeta$ derivative is $(2k-4t_0)/3<0$. Its value is at most
\[
 B(a)-\frac23(t_0-k/2)^2
 \le\frac{9550k-881}{3750}<\frac{41}{50}.
\]
The last comparison follows from $\sqrt2<577/408$, since
$577^2-2\cdot408^2=1$. Exchange $\xi,\zeta$ for the other
restriction. Finally $D\le\Phi-P<2k-41/50<3/350$, using
$\sqrt2<99/70$. This proves \eqref{eq:horizontal-short-data}.
In particular the positive diagonal tent is entirely inside $J$.

Put $\theta=\pi/8$, $c=\cos\theta$, $s=\sin\theta$,
$K=1/k$ and $K_0=Kc^2=(4+3\sqrt2)/4$.
The actual near-top support deficits are
\[
 \delta_L=c+s\zeta-h_U(5\pi/8),\qquad
 \delta_R=c+s\xi-h_U(3\pi/8).
\]
They are nonnegative and
$\delta_L\le s(\zeta-a/2)$, by the point $(-a/2,1)$.
Lowering the corresponding support line of $A_0$ removes the
triangle at $(-\zeta,1)$ with horizontal sides
$\delta_L/s$ and $\delta_L/(c-s)$, of area $K\delta_L^2$.
It lies in the left exterior wing. Indeed, if its second side
exceeded $a-\zeta$, the smaller triangle with deficit
$\delta_0=(c-s)(a-\zeta)$ would already give
\[
 D\ge K\delta_0^2=(a-\zeta)^2/\sqrt2
 >(13/100)^2/\sqrt2>3/350,
\]
contrary to \eqref{eq:horizontal-short-data}.
Reflection gives a disjoint right triangle. Consequently
\begin{equation}\label{eq:horizontal-short-exterior}
 D\ge K(\delta_L^2+\delta_R^2).
\end{equation}

Set $\eta=2/25$ and $E_L=(\eta-\delta_L/c)_+$,
$E_R=(\eta-\delta_R/c)_+$.
At the left zero $x_0=k-\zeta$ of the diagonal tent, the second
wall at $\theta$ has value
$1-1/c+k^2-\delta_L/c>E_L$ whenever $E_L>0$.
Thus it majorizes $E_L+k(x-x_0)$. A point attaining the diagonal
support has $x+y=1+\xi$, $y\le1$ and hence $x\ge\xi$;
it follows that $h_U(\pi/8)\ge s+c\xi$.
This first-wall lower bound majorizes the same line up to
$x_R=x_0+E_L/(1-k)$, because the comparison at $x_R$ is
\[
 2m\ge1/c+(1+\sqrt2)E_L.
\]
It holds by $m>16/25$, $E_L\le\eta$ and
$1/c+(1+\sqrt2)\eta<32/25$.
Also
\[
 x_R-d=k-m+E_L/(1-k)<0,
\]
since $k+\eta/(1-k)=(26k+3)/25<16/25$.
The added niche region is a triangle on the left of $d$, with
slopes $k$ and $1-k$, of area $K_0E_L^2$.
Its peak $x_0$ lies inside $J$, as
$x_0+a/2>k-8/25>0$.
Only its outer tip can be clipped, by length at most
\[
 L_0=8/25-k+\eta/k=(12-23k)/25>0.
\]
The lost area is at most $kL_0^2/2$.
The reflected construction at $3\pi/8$ lies right of $d$,
so the two payments are disjoint and
\begin{equation}\label{eq:horizontal-short-niche}
 \int_Jn\ge\int_Jn_{45}+K_0(E_L^2+E_R^2)-kL_0^2.
\end{equation}
When an $E$ vanishes its asserted area bound follows simply from
nonnegativity. The exact trigonometric comparisons above use
$c^2=(2+\sqrt2)/4$: squaring positive sides reduces
$1-1/c+k^2>2/25$ to $\sqrt2<1777/1249$, and
$1/c+(1+\sqrt2)2/25<32/25$ to $\sqrt2>231/167$.
These follow respectively from $\sqrt2<99/70$ and
$\sqrt2>24/17$.

Since $Kc^2=K_0$, for every $\delta\ge0$,
\[
 K\delta^2+K_0(\eta-\delta/c)_+^2\ge K_0\eta^2/2.
\]
Combining the two exterior and two niche payments gives
\[
 P_L(U)\le2k-K_0\eta^2+kL_0^2
 =\frac{5140k-1617}{625}
 <\frac{3587}{4375}<\frac{41}{50}.
\]
Here $k^2=1-2k$, $k<29/70$, and the last gap is $1/8750$.
This contradicts the supposition. Finally $4y^3+3y-1$ is
negative at $y=297/1000$, so $Y>297/1000$. The lower bound
$\arctan Y\ge Y-Y^3/3$ gives
\[
 \Mc/2>
 \bigl(1+4(297/1000)^2+297/1000-(297/1000)^3/3\bigr)/2
 >41/50.
\]
\end{proof}

\subsection{The top face and transfer between quarters}

For a full maximum the endpoint law gives, with $n_\pm=n(\pm C)$,
\begin{equation}\label{eq:horizontal-pressure}
 \begin{gathered}
 e_R=\frac{2+3n_+-n_-}{4},\qquad
 e_L=\frac{2+3n_--n_+}{4},\\
 E=e_R+e_L=1+(n_++n_-)/2.
 \end{gathered}
\end{equation}
Both end heights are positive.

\begin{lemma}\label{lem:horizontal-top}
If a horizontal canonical full maximum has $C\ge1/2$, its top
face is exactly $[-C,C]\times\{1\}$.
\end{lemma}
\begin{proof}
Write the top face $[a,b]\times\{1\}$; it contains $J$.
As $t\downarrow0$, $F(t)\to2C$ and $G(t)=1-at+o(t)$.
For $a<x<2C-1$ the first wall tends to $+\infty$, and
the second is $(x-a)t+o(t)>0$. Thus $n(x)>0$ there.
Reflection gives positivity on $(-2C+1,b)$.
For $C\ge1/2$ these intervals overlap or touch and cover the
interior of $J$ except possibly one point.
The positive-barrier projection bound of
\cref{cor:variation-projection} therefore gives
\[
 2C\le2C-(b-a-2C).
\]
Since $b-a\ge2C$, equality holds and the endpoints are $-C,C$.
Only lower semicontinuity on compact positive-barrier intervals
is used in this projection bound.
\end{proof}

The endpoint support values now give
\begin{equation}\label{eq:horizontal-corner-ends}
 p(0)=e_R,\quad q(0)=3C-1,\quad
 p(L)=1-3C,\quad q(L)=-e_L.
\end{equation}
On $\{q>0\}$ the curvature and shadow laws imply $p'\le-1/2$;
on $\{p<0\}$ they imply $q'\le-1/2$.
Indeed, for $0<q\le1$ the bound $u\le(1+q)/2$ gives
$p'\le-(1+q)/2$; for $q>1$, $u\le q$ gives $p'\le-1$.
For $-1\le p<0$, $v\le(1-p)/2$ gives
$q'\le(p-1)/2$; for $p<-1$, $v\le-p$ gives $q'\le-1$.
The positive-positive sector also has
\[
 p'=-1-q\le-1,\qquad q'\le p-1/2.
\]
Consequently any positive-$q$ component beginning at $t_A>0$
has $q\le1/8$: until $p$ becomes negative,
$p(t_A+s)\le1-s$ and the increase of $q$ is at most
$\int_0^{1/2}(1/2-s)\,ds=1/8$; after that crossing $q$
decreases. At a crossing $p=0$ with $q>0$, $p'<0$,
so the crossing on such a component is unique.

\begin{lemma}[Endpoint energy]\label{lem:horizontal-energy}
If $9C^2\le4+e_R(1-e_R)$, then $u\le1$ almost everywhere.
The reflected criterion with $e_L$ implies $v\le1$.
\end{lemma}
\begin{proof}
On the initial positive-positive sector,
\[
 \mathcal E=(p-1/2)^2+(q+1)^2,\qquad
 \mathcal E'=2(q+1)(v-1/2)\le0.
\]
If $p$ crosses zero on the initial positive-$q$ component, at
that crossing
$(q+1)^2\le9C^2-e_R(1-e_R)\le4$.
Thereafter $q$ decreases on that component, and later positive
components have $q\le1/8$. If the initial component ends before
the crossing, its positive-positive part has $u=0$.
Thus whenever $q>1$, one has $u=0$; elsewhere $-1\le q\le1$
and $u\le\kappa(q)\le1$.
\end{proof}

\begin{lemma}[One controlled quarter controls both]
\label{lem:horizontal-transfer}
For a horizontal canonical full maximum with $C>1/2$, either
$u\le1$ or $v\le1$ almost everywhere implies both inequalities.
\end{lemma}
\begin{proof}
Reflect so that $v\le1$, and suppose $u>1$ on a set of positive
measure. Such a point has $q>1$ and belongs to the initial
positive-$q$ component after its $p=0$ crossing.
This component ends at $\tau<L$. Let $b_0$ be its unique
subsequent $q=1$ crossing; then $p<0<q<1$ on $(b_0,\tau)$.
Every future positive component has $q\le1/8$, so $u\le1$
on $(b_0,L)$.

For the inner corner $c=(F-1)u(t)+(G-1)u^\perp(t)$ put
\[
 B=c+p u^\perp(t),\qquad D=c-q u(t).
\]
Differentiation gives
$B'=(u-1)u^\perp(t)$ and $D'=(1-v)u(t)$.
The top endpoints give $D(0)=(-C,0)$ and $B(L)=(C,0)$.
Thus $D_x,D_y$ are globally nondecreasing, and $B_x$ is
nondecreasing after $b_0$. On $(b_0,\tau)$,
\begin{equation}\label{eq:horizontal-visible-window}
 -C\le D_x<c_x<B_x\le C,\qquad c_y=D_y+q\sin t>0.
\end{equation}
The $D$ graph has positive height wherever its arc density $1-v$
is nonzero: the derivative of $D_y$ vanishes almost everywhere
on its zero level set. Its contributing points lie strictly inside $J$.

The wall parameter derivatives are
\[
 \partial_t a_t(x)=\frac{x-B_x(t)}{\sin^2t},\qquad
 \partial_t b_t(x)=\frac{x-D_x(t)}{\cos^2t}.
\]
Global monotonicity of $D_x$ makes $D(t)$ a maximum of the
entire second-wall family, with positive companion gap $q/\sin t$.
Unless $D$ is constant on a parameter interval, which has zero
arc contribution, its active source parameter is unique.
At $c(t)$ every earlier second wall is strictly lower, because
$c_x(t)-D_x(s)\ge q(t)\cos t>0$ for $s<t$;
every later first wall is strictly lower, because
$c_x(t)-B_x(s)\le p(t)\sin t<0$ for $s>t$.
The corner likewise has unique global source parameter and
strictly decreasing abscissa.
For $K=[t_0,t_1]\subset(b_0,\tau)$ the two graph images are
disjoint, since $D_x(t_1)<c_x(t_1)$.

We can therefore add their finite source fluxes. On positive-height
compact subarcs, active finite source parameters converge to the
unique limiting parameter: any contrary subsequence would supply
a different maximizing parameter by uniform wall convergence.
The angles stay away from the axes, so graph slopes are bounded.
With increasing abscissa the finite vector flux is
\[
 (\dd x,\dd y)=-u^\perp(t)\,\dd\nu_i^f
                         +u(t)\,\dd\nu_i^g.
\]
Its first component is Lebesgue measure, and its second converges
by uniform roof convergence and the common slope bound.
Solving the two-direction system gives
$\dd\nu^g=\cos t\,\dd x+\sin t\,\dd y$.
The $D$ graph therefore contributes $(1-v)\dd t$ to this source,
and the reverse-oriented corner contributes $-p\dd t$.
Exhaustion removes constant and zero-height pieces.
Source equality from \cref{lem:variation-sources} yields
\[
 2v\ge1-p\quad\hbox{almost everywhere on }(b_0,\tau).
\]
Since $v\le1$, $p\ge-1$ there. The upper curvature bound gives
the reverse inequality $v\le(1-p)/2$, hence equality.
Set $x=1-p$, $y=1+q$. Then
\[
 x'=1+q-u\ge y/2,\qquad y'=-x/2,\qquad
 (x^2+y^2)'\ge0.
\]
At $b_0$, $y=2$ and $x>1$, so this energy exceeds five;
at $\tau$, $y=1$ and $x\le2$, so it is at most five.
This contradiction proves the transfer.
\end{proof}

\begin{lemma}[Box bound and the narrow-width criterion]
\label{lem:horizontal-box}
For $C<1$, $n(\pm C)\le H(C)=1-\sqrt{1-C^2}$.
Every horizontal canonical full maximum with $C\le C_H$
has $u,v\le1$ almost everywhere.
\end{lemma}
\begin{proof}
The box gives $a_t(C)\le1+(C\cos t-1)/\sin t$.
Its maximum is $H(C)$, at $\cos t=C$. Reflection treats $-C$.
Put $x=n(C)$, $y=n(-C)$. By \eqref{eq:horizontal-pressure},
$e_R-1/2=(3x-y)/4$ and $e_L-1/2=(3y-x)/4$.
At least one of these absolute values is at most $H(C)/2$.
Indeed, if $x\ge y$ and $3y\ge x$, then
$0\le(3y-x)/4\le H/2$; if $3y<x$ its absolute value
is at most $H/4$. Exchange $x,y$ otherwise.
Thus at least one $e(1-e)\ge(1-H^2)/4$.
The energy criterion applies when
\[
 9C^2\le4+(1-H(C)^2)/4,
 \quad\hbox{equivalently}\quad
 35C^2\le15+2\sqrt{1-C^2}.
\]
The difference is strictly increasing in $C^2$, and its unique
zero is $(523+2\sqrt{701})/1225$.
Use \cref{lem:horizontal-energy,lem:horizontal-transfer}.
The impossible range $C\le1/2$ is already excluded by
\cref{lem:horizontal-short}.
\end{proof}

\subsection{Two strict wide-width exclusions}

Return to $a=2C$ and the diagonal data $\xi,\zeta,m$.
The endpoint points give $\xi\ge a+e_R-1$,
$\zeta\ge a+e_L-1$; since $e_R+e_L\ge1$,
\begin{equation}\label{eq:horizontal-diagonal-constraint}
 m\ge a-1/2.
\end{equation}
For $C\ge13/15$, this is to the right of the symmetric
unclipped maximum $C+k/2$ of $\Phi(a,m,m)$.
In the clipped region its $m$ derivative is $2(C-m)<0$.
Evaluating at $m=2C-1/2$ gives
\[
 P_L(U)\le
 \begin{cases}
 2C-1/4-(2C-1/2-k)^2,&C\le1/2+k,\\
 (1+2\sqrt2)C-3C^2-1/4,&C\ge1/2+k.
 \end{cases}
\]
Both branches decrease on the relevant ranges, agree at the
switch, and the first at $13/15$ is
\[
 \frac{67\sqrt2}{15}-\frac{2477}{450}
 <\frac{256}{315}<\frac{41}{50}.
\]
Thus any horizontal full maximum has $C<13/15$.

\begin{lemma}[Three-angle wide exclusion]\label{lem:horizontal-wide}
Every horizontal canonical full maximum has $C<4/5$.
\end{lemma}
\begin{proof}
It remains to consider $8/5\le a=2C\le26/15$.
The box bound gives $H(C)<2/3$, so the pressures imply
$e_R,e_L\ge1/3$ and $E\ge1$.
Put $w_R=a-\xi$, $w_L=a-\zeta$. Then
$0\le w_R,w_L\le2/3$ and $m\ge a-1/2\ge11/10$.
As before let $c=\cos(\pi/8)$, $s=\sin(\pi/8)$ and
\[
 \delta_L=c+s\zeta-h_U(5\pi/8),\quad
 \delta_R=c+s\xi-h_U(3\pi/8).
\]
Here $\delta_L\le s(\zeta-a/2)$.
On the flat part $x=-\zeta+z$, $0\le z\le\delta_L/s$,
the new support line cuts height $\delta_L/c-kz$ from $A_0$.
This triangle is entirely in the left wing and has area
$\delta_L^2/(2sc)=\sqrt2\delta_L^2$.
Reflection supplies a disjoint right triangle, giving
\begin{equation}\label{eq:horizontal-wide-exterior}
 D\ge\sqrt2(\delta_L^2+\delta_R^2).
\end{equation}

The actual second wall at $\pi/8$ is $b+kx$, where
$b=1-1/c+k\zeta-\delta_L/c$.
The endpoint $(a,e_R)$ bounds its first companion below by
$B-x/k$, $B=e_R+a/k-1/s$.
Define
\[
 x_L=(\xi-k-b)/\sqrt2,\quad
 x_*=(B-b)/(2\sqrt2),\quad
 Z_R=(B+b)/2-\xi+k.
\]
If $Z_R>0$, on $[x_L,x_*]$ the lower tent $b+kx$
exceeds the diagonal roof $\xi-k-x$ by $\sqrt2(x-x_L)$.
The entire added half-triangle is inside $J$ and right of $d$.
For completeness the needed bounds are
\[
 x_L-d=km+2s-1+\delta_L/(\sqrt2c)>0,
\]
using $m\ge11/10$, $k>2/5$, $s>3/8$; and
\[
 x_*=\frac{2a+e_R+kw_L-(1+4s)+\delta_L/c}{2\sqrt2}<a/2.
\]
For the latter use $\delta_L/c\le k(a/2-w_L)$ and
$((3-\sqrt2)/2)a\le4s$, since the left side is
less than $104/75$ and the right side exceeds $3/2$.
Finally $x_*<\xi-k$: using $e_R\le1-w_R$, $w_R\le2/3$,
it suffices that
\[
 (2\sqrt2-1)2/3\le(12/5)k+4s-2\sqrt2k.
\]
The left side is below $26/21$ and the right side above
$3/2-48/245=639/490$, using $7/5<\sqrt2<10/7$,
$k<3/7$, $s>3/8$.
Thus the diagonal roof is positive on the half-triangle,
whose added area is $Z_R^2/(2\sqrt2)$.
For $Z_R\le0$ the zero lower bound needs no interval claim.

Reflection at $3\pi/8$ gives a left half-triangle of area
$(Z_L)_+^2/(2\sqrt2)$, disjoint from the first.
Writing $Z_R=Z_R^0-\delta_L/(2c)$ and
$Z_L=Z_L^0-\delta_R/(2c)$, their uncut sum is
\[
 S_0=Z_R^0+Z_L^0
 =E/2+a/k+(k-2)m+1-1/s-1/c+2k.
\]
Put $D_0=\sqrt2(\delta_L^2+\delta_R^2)$ and
$G_0'=((Z_R)_+^2+(Z_L)_+^2)/(2\sqrt2)$.
Cauchy--Schwarz gives
\[
 (S_0)_+\le\sqrt{4\sqrt2G_0'}+
              \sqrt{\sqrt2D_0/(4c^2)},
 \qquad (S_0)_+^2\le8(D_0+G_0').
\]
The coefficient $4\sqrt2+\sqrt2/(4c^2)$ is below eight.
Hence the exterior and niche payments give
\[
 P_L(U)\le\Phi(a,m,m)-(S_0)_+^2/8.
\]
Using $E\ge1$, $1/k>12/5$, $k-2>-8/5$ and
$3/2+2k-4c>-11/8$, we may replace $S_0$ by
$S=(12/5)a-(8/5)m-11/8$.
The last constant inequality follows by squaring
$4c<7/8+2\sqrt2$; the squared difference is
$49/64-\sqrt2/2>0$.

The resulting function $\mathcal G=\Phi(a,m,m)-(S_+)^2/8$
is jointly concave. Before and after clipping, the Hessians of
$\Phi(a,m,m)$ are respectively
$\left(\begin{smallmatrix}-2&2\\2&-4\end{smallmatrix}\right)$
and $\left(\begin{smallmatrix}-3/2&1\\1&-2\end{smallmatrix}\right)$;
their first derivatives match at the boundary.
At $(a_0,m_0)=(8/5,11/10)$ the clipping vanishes, $S=141/200$,
and
\[
 \mathcal G_a=-423/1000,\qquad
 \mathcal G_m=2k-459/500<0.
\]
Every allowed displacement is
$(a-a_0)(1,1)+(m-a+1/2)(0,1)$ with nonnegative coefficients.
Both directional derivatives are negative. The supporting-plane
bound therefore gives
\[
 P_L(U)\le\mathcal G(a_0,m_0)
 =\frac{21}5k-\frac{295081}{320000}
 <\frac{261719}{320000}<\frac{41}{50}<\Mc/2,
\]
a contradiction.
\end{proof}

\subsection{The leakage estimate and its exact rational intervals}

The remaining exclusion uses the ordinary-sofa bound with
$G_0=22199/10000$. We first check its geometric applicability.
For $1/2\le C\le13/15$, every inner corner has $-C\le c_x\le C$.
Indeed, with $s=\sin t$, $c=\cos t$, the box and left top point give
$F\le2Cc+s$, $G\ge Cs+c$, so
$c_x\le C(2-3s^2)-c+s$.
The required inequality $C(1-3s^2)\le c-s$ follows in three ranges.
For $s^2\le1/3$, use $c\ge1-3s^2/5$ and $C\le13/15$;
the remaining margin is
$2s^2-s+2/15=2(s-1/4)^2+1/120>0$.
For $1/3\le s^2\le1/2$ the two sides have opposite signs.
For $s^2\ge1/2$, use $C\ge1/2$ and
\[
 2(s-c)=2(2s^2-1)/(s+c)\le4s^2-2\le3s^2-1.
\]
Reflection supplies the other bound.

Consequently $n$ is convex on each wing: there each quadrant
uses the same descending, or reflected ascending, wall.
The box gives $n(\pm2C)=0$ and zero niche outside $I$.
On $J$, \cref{cor:variation-green} gives $n\le E/2\le1=A$;
on each wing $A-n$ is concave and nonnegative at both endpoints.
Thus $n\le A$ everywhere, and \cref{lem:ordinary-survivor}
applies to the actual survivor.
Wing convexity yields
$N_{\rm out}\le C(n_++n_-)/2=C(E-1)$.
The wing-length identity and its chord bounds yield
\[
 \begin{aligned}
 P_L(U)=L_{\rm wing}/2
 &\ge\tfrac12\bigl(\sqrt{C^2+(1-e_R)^2}
                  +\sqrt{C^2+(1-e_L)^2}\bigr)\\
 &\ge\sqrt{C^2+(1-E/2)^2}.
 \end{aligned}
\]
Since the survivor area is $P_L(U)+2C-N_{\rm out}$,
the ordinary bound gives
\begin{equation}\label{eq:horizontal-ordinary-necessary}
 C(3-E)+\sqrt{C^2+(1-E/2)^2}\le G_0.
\end{equation}
If $n_\pm\le B\le1$, its left side decreases with $E\le1+B$,
so necessarily
\begin{equation}\label{eq:horizontal-ordinary-scalar}
 C(2-B)+\sqrt{C^2+(1-B)^2/4}\le G_0.
\end{equation}

Fix the right endpoint $e=e_R$.
Any $u>1$ lies after the initial $p=0$ crossing $t_A$ and before
the first subsequent $q=1$ crossing. The energy gives
\[
 \begin{gathered}
 q(t_A)+1\le\sqrt{9C^2-e(1-e)},\\
 \epsilon=(\sqrt{9C^2-e(1-e)}-2)_+,\qquad
 h=\sqrt{1+4\epsilon}-1.
 \end{gathered}
\]
While $q>1$, $p'\le-1$, hence $p(t_A+\tau)\le-\tau$.
For $0\le\tau\le1$, $q'\le-(1+\tau)/2$:
if $p\ge-1$, use $v\le(1-p)/2$; if $p<-1$,
use $v\le-p$, giving the stronger $q'\le-1$.
Therefore
\begin{equation}\label{eq:horizontal-leakage-envelope}
 (u(t_A+\tau)-1)_+
 \le(\epsilon-\tau/2-\tau^2/4)_+.
\end{equation}
The endpoint box bound and pressure law give
$|e-1/2|\le3H(C)/4$ and
\[
 \epsilon\le
 (\sqrt{9C^2-1/4+9H(C)^2/16}-2)_+<2/5
 \quad(C\le4/5).
\]
Thus $h<1$ covers every possible excess interval.
Before the crossing, $q'\ge-1$ and $p'=-1-q\le-3C+t$,
so $p(t)\le e-3Ct+t^2/2$. Its first root gives
\begin{equation}\label{eq:horizontal-crossing-time}
 t_A\le T(C,e)=3C-\sqrt{9C^2-2e}.
\end{equation}
If the component ends before a crossing, its leakage is zero.

The top endpoint data $F(L)=1$, $F'(L)=-C$ give the exact
backward solution
\[
 a_t(C)=\frac1{\sin t}\int_t^L\sin(s-t)(u(s)-1)\,\dd s.
\]
The box makes this wall nonpositive unless
\[
 t>t_{\rm box}(C)=2\arctan\frac{1-C}{1+C},\qquad
 \sin t_{\rm box}=\frac{1-C^2}{1+C^2}.
\]
Discarding negative curvature differences and using
$\sin(s-t)\le s-t$ in the remaining positive integral yields
\begin{equation}\label{eq:horizontal-leakage-kernel}
 n(C)\le
 \frac{A_0(t_A-t_{\rm box})_++A_1}{\sin t_{\rm box}},
 \qquad A_0=h^2/4+h^3/6,\quad A_1=h^3/12+h^4/16.
\end{equation}
These are the integrals of the envelope and of $\tau$ times
the envelope in \eqref{eq:horizontal-leakage-envelope}; use
$\epsilon=(h^2+2h)/4$.
Even when $t_A<t_{\rm box}$, the bound is valid because
$s-t\le(t_A-t_{\rm box})_++\tau$.
Reflection gives the same estimate at $-C$.

\begin{lemma}[Exact leakage exclusion]\label{lem:horizontal-leakage}
No horizontal canonical full maximum has
$8571/12500\le C\le4/5$.
\end{lemma}
\begin{proof}
Every terminating decimal in the following table is an exact
rational. A row gives an interval $[a,b]$ and witnesses
$H_*,h_*,T_*,t_*,B_*$.
\begin{center}\small
\begin{tabular}{rrrrrrr}
\toprule
$a$&$b$&$H_*$&$h_*$&$T_*$&$t_*$&$B_*$\\
\midrule
0.68568&0.686&0.2724&0.0135&0.377&0.368&$1/400000$\\
0.686&0.69&0.2762&0.038&0.3784&0.3627&$1/30000$\\
0.69&0.7&0.2859&0.0971&0.3801&0.3492&$1/2000$\\
0.7&0.725&0.3113&0.2323&0.3845&0.3161&$3/400$\\
0.725&0.75&0.3386&0.3542&0.3799&0.2837&$31/1000$\\
0.75&0.775&0.3681&0.4662&0.3765&0.2521&$83/1000$\\
0.775&0.8&0.4&0.5703&0.3742&0.2213&$9/50$\\
\bottomrule
\end{tabular}
\end{center}
Put $z=(1-b)/(1+b)$, $\epsilon_*=(h_*^2+2h_*)/4$;
evaluate $A_0,A_1$ at $h_*$.
The following five rational residuals are nonnegative:
\begin{align*}
 \Delta_1&=1-b^2-(1-H_*)^2,\\
 \Delta_2&=(2+\epsilon_*)^2-9b^2+1/4-9H_*^2/16,\\
 \Delta_3&=6aT_*-T_*^2-1-3H_*/2,\\
 \Delta_4&=2(z-z^3/3)-t_*,\\
 \Delta_5&=B_*\frac{1-b^2}{1+b^2}
                   -A_0(T_*-t_*)_+-A_1.
\end{align*}
Also $T_*<3a$. Thus $H(C)\le H_*$, $h\le h_*$,
$t_A\le T_*$ by squaring \eqref{eq:horizontal-crossing-time},
and $t_{\rm box}(C)\ge t_*$ by
$\arctan z\ge z-z^3/3$. The last residual and
\eqref{eq:horizontal-leakage-kernel} give $n_\pm\le B_*$.
For an explicit record of the rational sign checks, the entries
below are the numerators after reducing each residual to lowest
terms; every omitted denominator is positive. The last column
uses $\Delta_6=a^2+(1-B_*)^2/4-R_*^2$,
$R_*=G_0-a(2-B_*)$.
\begin{center}\scriptsize
\begin{tabular}{rrrrrr}
\toprule
$\Delta_1$&$\Delta_2$&$\Delta_3$&$\Delta_4$&$\Delta_5$&$\Delta_6$\\
\midrule
7&32240555441&6979&38805734&1050817008759&3319291602200459\\
339&7725821&49&7131271&291483040799&120460934617\\
6119&38697993591281&28799&951&167988627540997&793536099\\
6731&37513157356241&439&389153&46533834362383573&2769577\\
1251&3244477366481&34099&727&7903542862229&1696091\\
7739&5192642710721&1391&226369&11606973036545727&64385451\\
0&21664728155681&109&169&1866283140377871&709979\\
\bottomrule
\end{tabular}
\end{center}
The respective $R_*$ are
\[
 \frac{4242708571}{5000000000},\quad
 \frac{12718843}{15000000},\quad\frac{168049}{200000},\quad
 \frac{16503}{20000},\quad\frac{6339}{8000},\quad
 \frac{15643}{20000},\quad\frac{4047}{5000};
\]
all are positive. Therefore $\Delta_6>0$ proves
$a(2-B_*)+\sqrt{a^2+(1-B_*)^2/4}>G_0$.
The left side increases with $C$ at fixed $B_*$, contradicting
\eqref{eq:horizontal-ordinary-scalar} throughout each row.
The intervals cover the claimed range.
\end{proof}

\begin{proof}[Proof of \cref{lem:full-horizontal-unit}]
The short and wide exclusions give $1/2<C<4/5$, and the
leakage exclusion gives $C<8571/12500$.
This upper bound is below $C_H$. Indeed, with
$z_0=(8571/12500)^2$, $35z_0-15>0$ and
\[
 4(1-z_0)-(35z_0-15)^2
 =\frac{878611101631}{976562500000000}>0.
\]
Hence $35z_0-15-2\sqrt{1-z_0}<0$.
\Cref{lem:horizontal-top,lem:horizontal-box} now give the
exact top face and both unit curvature bounds.

For the final sign statement, $B_x,D_x$ are nondecreasing and
\[
 B_x-D_x=q\cos t-p\sin t\ge4C-2>0,
\]
because $B_x\ge B_x(0)=2C-1$ and
$D_x\le D_x(L)=1-2C$.
At any zero of $p$, this strict inequality gives $q>0$,
so $p'=u-1-q<0$ nearby; a return through zero is impossible.
At any zero of $q$, it gives $p<0$, so $q'=v-1+p<0$
nearby, likewise excluding a return.
The endpoint signs in \eqref{eq:horizontal-corner-ends}
give one crossing each. They cannot coincide, and at the $p$
crossing $q>0$, so their order is $\beta<\gamma$.
\end{proof}
